% ---- Acronyms & Glossary entries ----
% Define your acronyms and glossary terms here.
% Usage in text:
%   \acrfull{key}   → full form + acronym (e.g. "واجهة الدماغ–الحاسوب (BCI)")
%   \acrshort{key}  → acronym only (e.g. "BCI")
%   \acrlong{key}   → long form only
%   \gls{key}       → first use: full; subsequent: short
%
% Format: \newacronym{key}{acronym}{long form (Arabic) - English}

\renewcommand*{\acronymname}{الاختصارات}
\renewcommand*{\glossaryname}{المصطلحات}

% ---- Example acronyms (replace with your own) ----
\newacronym{bci}{\LR{BCI}}{\textarabic{ترابط الدماغ–الحاسوب} - \LR{Brain-Computer Interface}}
\newacronym{eeg}{\LR{EEG}}{\textarabic{تخطيط الدماغ الكهربائي} - \LR{Electroencephalography}}
\newacronym{ai}{\LR{AI}}{\textarabic{الذكاء الاصطناعي} - \LR{Artificial Intelligence}}
\newacronym{ml}{\LR{ML}}{\textarabic{تعلُّم الآلة} - \LR{Machine Learning}}

% ---- Example glossary entries (not acronyms) ----
\newglossaryentry{neuron}{
  name=\LR{Neuron},
  description={العَصَبُون: الوحدة الوظيفية الأساسية للجهاز العصبي، مسؤول عن نقل الإشارات الكهربائية بين الخلايا}
}

\newglossaryentry{neuroplasticity}{
  name=\LR{Neuroplasticity},
  description={اللُّدُونَةُ العَصَبِيَّة: قدرة الجهاز العصبي على التكيّف وإعادة تنظيم نفسه بنيويًا ووظيفيًا نتيجة الخبرة أو الإصابة أو التعلم}
}
